% ============================================================
% tables_memoria.tex
%
% Tabelas de uso de memória para algoritmos PQC (TCC).
% Inclua no documento com: % ============================================================
% tables_memoria.tex
%
% Tabelas de uso de memória para algoritmos PQC (TCC).
% Inclua no documento com: % ============================================================
% tables_memoria.tex
%
% Tabelas de uso de memória para algoritmos PQC (TCC).
% Inclua no documento com: % ============================================================
% tables_memoria.tex
%
% Tabelas de uso de memória para algoritmos PQC (TCC).
% Inclua no documento com: \input{tables_memoria}
%
% Pacotes necessários:
%   \usepackage{booktabs}
%   \usepackage{multirow}
%   \usepackage[utf8]{inputenc}
%
% Metodologia de medição:
%   - Stack  : stack painting via pthread (4 MB de buffer sentinela)
%   - Heap   : Valgrind --tool=massif, pico de mem_heap_B
% ============================================================


% ============================================================
% Tabela 1: Stack — KEMs
% ============================================================
\begin{table}[ht]
\centering
\caption{Uso de memória \textit{stack} por mecanismo de encapsulamento de
chave (KEM) e operação, em bytes. Medido via \textit{stack painting}
com \texttt{pthread} de 4\,MB e compilação \texttt{-O0}.}
\label{tab:stack_kem}
\begin{tabular}{lrrr}
\toprule
\textbf{Algoritmo} & \textbf{Keygen (B)} & \textbf{Encap (B)} & \textbf{Decap (B)} \\
\midrule
\multicolumn{4}{l}{\textit{ML-KEM (FIPS 203)}} \\[2pt]
\quad ML-KEM-512  &  13\,520 &  16\,160 &  16\,960 \\
\quad ML-KEM-768  &  17\,648 &  20\,784 &  21\,904 \\
\quad ML-KEM-1024 &  22\,768 &  26\,448 &  28\,048 \\
\midrule
\multicolumn{4}{l}{\textit{HQC}} \\[2pt]
\quad HQC-128 &  38\,576 &  56\,448 &  60\,896 \\
\quad HQC-192 &  70\,688 & 106\,752 & 115\,760 \\
\quad HQC-256 & 108\,768 & 166\,592 & 181\,040 \\
\bottomrule
\end{tabular}
\end{table}


% ============================================================
% Tabela 2: Stack — Assinaturas Digitais
% ============================================================
\begin{table}[ht]
\centering
\caption{Uso de memória \textit{stack} por esquema de assinatura digital
e operação, em bytes. Medido via \textit{stack painting}
com \texttt{pthread} de 4\,MB e compilação \texttt{-O0}.
Dados para SLH-DSA-SHAKE não disponíveis para esta métrica
(medição pendente com variantes SHAKE).}
\label{tab:stack_sig}
\begin{tabular}{lrrr}
\toprule
\textbf{Algoritmo} & \textbf{Keygen (B)} & \textbf{Sign (B)} & \textbf{Verify (B)} \\
\midrule
\multicolumn{4}{l}{\textit{ML-DSA (FIPS 204)}} \\[2pt]
\quad ML-DSA-44 &  28\,304 &  53\,456 &  25\,968 \\
\quad ML-DSA-65 &  33\,456 &  83\,312 &  29\,072 \\
\quad ML-DSA-87 &  41\,648 & 117\,072 &  35\,472 \\
\midrule
\multicolumn{4}{l}{\textit{FN-DSA (FIPS 206)}} \\[2pt]
\quad FN-DSA-512  &  23\,608 &  49\,584 &   9\,064 \\
\quad FN-DSA-1024 &  40\,504 &  90\,736 &  13\,160 \\
\midrule
\multicolumn{4}{l}{\textit{SLH-DSA-SHAKE (FIPS 205)}} \\[2pt]
\quad SLH-DSA-SHAKE-128s & --- & --- & --- \\
\quad SLH-DSA-SHAKE-128f & --- & --- & --- \\
\quad SLH-DSA-SHAKE-192s & --- & --- & --- \\
\quad SLH-DSA-SHAKE-192f & --- & --- & --- \\
\quad SLH-DSA-SHAKE-256s & --- & --- & --- \\
\quad SLH-DSA-SHAKE-256f & --- & --- & --- \\
\bottomrule
\end{tabular}
\end{table}


% ============================================================
% Tabela 3: Pico de Heap — KEMs (Valgrind massif)
% ============================================================
\begin{table}[ht]
\centering
\caption{Pico de uso de memória \textit{heap} por mecanismo de
encapsulamento de chave (KEM) e operação, em KB.
Medido com \texttt{valgrind --tool=massif}.
Os valores representam o máximo de \texttt{mem\_heap\_B}
ao longo da execução de cada operação isolada.}
\label{tab:heap_kem}
\begin{tabular}{lrrr}
\toprule
\textbf{Algoritmo} & \textbf{Keygen (KB)} & \textbf{Encap (KB)} & \textbf{Decap (KB)} \\
\midrule
\multicolumn{4}{l}{\textit{ML-KEM (FIPS 203)}} \\[2pt]
\quad ML-KEM-512  & 204.97 & 204.97 & 204.97 \\
\quad ML-KEM-768  & 207.17 & 207.17 & 207.17 \\
\quad ML-KEM-1024 & 208.00 & 208.17 & 208.17 \\
\midrule
\multicolumn{4}{l}{\textit{HQC}} \\[2pt]
\quad HQC-128 & 211.07 & 211.07 & 211.07 \\
\quad HQC-192 & 219.74 & 219.74 & 219.74 \\
\quad HQC-256 & 230.60 & 230.60 & 230.60 \\
\bottomrule
\end{tabular}
\end{table}


% ============================================================
% Tabela 4: Pico de Heap — Assinaturas Digitais (Valgrind massif)
% ============================================================
\begin{table}[ht]
\centering
\caption{Pico de uso de memória \textit{heap} por esquema de assinatura
digital e operação, em KB.
Medido com \texttt{valgrind --tool=massif}.
Os valores representam o máximo de \texttt{mem\_heap\_B}
ao longo da execução de cada operação isolada.}
\label{tab:heap_sig}
\begin{tabular}{lrrr}
\toprule
\textbf{Algoritmo} & \textbf{Keygen (KB)} & \textbf{Sign (KB)} & \textbf{Verify (KB)} \\
\midrule
\multicolumn{4}{l}{\textit{ML-DSA (FIPS 204)}} \\[2pt]
\quad ML-DSA-44 & 207.92 & 208.90 & 208.90 \\
\quad ML-DSA-65 & 211.02 & 211.02 & 211.83 \\
\quad ML-DSA-87 & 214.37 & 214.59 & 214.59 \\
\midrule
\multicolumn{4}{l}{\textit{FN-DSA (FIPS 206)}} \\[2pt]
\quad FN-DSA-512  & 204.64 & 204.81 & 204.81 \\
\quad FN-DSA-1024 & 207.16 & 207.38 & 207.38 \\
\midrule
\multicolumn{4}{l}{\textit{SLH-DSA-SHAKE (FIPS 205)}} \\[2pt]
\quad SLH-DSA-SHAKE-128s & 210.31 & 210.31 & 210.31 \\
\quad SLH-DSA-SHAKE-128f & 219.32 & 219.32 & 219.32 \\
\quad SLH-DSA-SHAKE-192s & 218.52 & 218.52 & 218.52 \\
\quad SLH-DSA-SHAKE-192f & 237.51 & 237.51 & 237.51 \\
\quad SLH-DSA-SHAKE-256s & 231.82 & 231.82 & 231.82 \\
\quad SLH-DSA-SHAKE-256f & 251.42 & 251.42 & 251.42 \\
\bottomrule
\end{tabular}
\end{table}

%
% Pacotes necessários:
%   \usepackage{booktabs}
%   \usepackage{multirow}
%   \usepackage[utf8]{inputenc}
%
% Metodologia de medição:
%   - Stack  : stack painting via pthread (4 MB de buffer sentinela)
%   - Heap   : Valgrind --tool=massif, pico de mem_heap_B
% ============================================================


% ============================================================
% Tabela 1: Stack — KEMs
% ============================================================
\begin{table}[ht]
\centering
\caption{Uso de memória \textit{stack} por mecanismo de encapsulamento de
chave (KEM) e operação, em bytes. Medido via \textit{stack painting}
com \texttt{pthread} de 4\,MB e compilação \texttt{-O0}.}
\label{tab:stack_kem}
\begin{tabular}{lrrr}
\toprule
\textbf{Algoritmo} & \textbf{Keygen (B)} & \textbf{Encap (B)} & \textbf{Decap (B)} \\
\midrule
\multicolumn{4}{l}{\textit{ML-KEM (FIPS 203)}} \\[2pt]
\quad ML-KEM-512  &  13\,520 &  16\,160 &  16\,960 \\
\quad ML-KEM-768  &  17\,648 &  20\,784 &  21\,904 \\
\quad ML-KEM-1024 &  22\,768 &  26\,448 &  28\,048 \\
\midrule
\multicolumn{4}{l}{\textit{HQC}} \\[2pt]
\quad HQC-128 &  38\,576 &  56\,448 &  60\,896 \\
\quad HQC-192 &  70\,688 & 106\,752 & 115\,760 \\
\quad HQC-256 & 108\,768 & 166\,592 & 181\,040 \\
\bottomrule
\end{tabular}
\end{table}


% ============================================================
% Tabela 2: Stack — Assinaturas Digitais
% ============================================================
\begin{table}[ht]
\centering
\caption{Uso de memória \textit{stack} por esquema de assinatura digital
e operação, em bytes. Medido via \textit{stack painting}
com \texttt{pthread} de 4\,MB e compilação \texttt{-O0}.
Dados para SLH-DSA-SHAKE não disponíveis para esta métrica
(medição pendente com variantes SHAKE).}
\label{tab:stack_sig}
\begin{tabular}{lrrr}
\toprule
\textbf{Algoritmo} & \textbf{Keygen (B)} & \textbf{Sign (B)} & \textbf{Verify (B)} \\
\midrule
\multicolumn{4}{l}{\textit{ML-DSA (FIPS 204)}} \\[2pt]
\quad ML-DSA-44 &  28\,304 &  53\,456 &  25\,968 \\
\quad ML-DSA-65 &  33\,456 &  83\,312 &  29\,072 \\
\quad ML-DSA-87 &  41\,648 & 117\,072 &  35\,472 \\
\midrule
\multicolumn{4}{l}{\textit{FN-DSA (FIPS 206)}} \\[2pt]
\quad FN-DSA-512  &  23\,608 &  49\,584 &   9\,064 \\
\quad FN-DSA-1024 &  40\,504 &  90\,736 &  13\,160 \\
\midrule
\multicolumn{4}{l}{\textit{SLH-DSA-SHAKE (FIPS 205)}} \\[2pt]
\quad SLH-DSA-SHAKE-128s & --- & --- & --- \\
\quad SLH-DSA-SHAKE-128f & --- & --- & --- \\
\quad SLH-DSA-SHAKE-192s & --- & --- & --- \\
\quad SLH-DSA-SHAKE-192f & --- & --- & --- \\
\quad SLH-DSA-SHAKE-256s & --- & --- & --- \\
\quad SLH-DSA-SHAKE-256f & --- & --- & --- \\
\bottomrule
\end{tabular}
\end{table}


% ============================================================
% Tabela 3: Pico de Heap — KEMs (Valgrind massif)
% ============================================================
\begin{table}[ht]
\centering
\caption{Pico de uso de memória \textit{heap} por mecanismo de
encapsulamento de chave (KEM) e operação, em KB.
Medido com \texttt{valgrind --tool=massif}.
Os valores representam o máximo de \texttt{mem\_heap\_B}
ao longo da execução de cada operação isolada.}
\label{tab:heap_kem}
\begin{tabular}{lrrr}
\toprule
\textbf{Algoritmo} & \textbf{Keygen (KB)} & \textbf{Encap (KB)} & \textbf{Decap (KB)} \\
\midrule
\multicolumn{4}{l}{\textit{ML-KEM (FIPS 203)}} \\[2pt]
\quad ML-KEM-512  & 204.97 & 204.97 & 204.97 \\
\quad ML-KEM-768  & 207.17 & 207.17 & 207.17 \\
\quad ML-KEM-1024 & 208.00 & 208.17 & 208.17 \\
\midrule
\multicolumn{4}{l}{\textit{HQC}} \\[2pt]
\quad HQC-128 & 211.07 & 211.07 & 211.07 \\
\quad HQC-192 & 219.74 & 219.74 & 219.74 \\
\quad HQC-256 & 230.60 & 230.60 & 230.60 \\
\bottomrule
\end{tabular}
\end{table}


% ============================================================
% Tabela 4: Pico de Heap — Assinaturas Digitais (Valgrind massif)
% ============================================================
\begin{table}[ht]
\centering
\caption{Pico de uso de memória \textit{heap} por esquema de assinatura
digital e operação, em KB.
Medido com \texttt{valgrind --tool=massif}.
Os valores representam o máximo de \texttt{mem\_heap\_B}
ao longo da execução de cada operação isolada.}
\label{tab:heap_sig}
\begin{tabular}{lrrr}
\toprule
\textbf{Algoritmo} & \textbf{Keygen (KB)} & \textbf{Sign (KB)} & \textbf{Verify (KB)} \\
\midrule
\multicolumn{4}{l}{\textit{ML-DSA (FIPS 204)}} \\[2pt]
\quad ML-DSA-44 & 207.92 & 208.90 & 208.90 \\
\quad ML-DSA-65 & 211.02 & 211.02 & 211.83 \\
\quad ML-DSA-87 & 214.37 & 214.59 & 214.59 \\
\midrule
\multicolumn{4}{l}{\textit{FN-DSA (FIPS 206)}} \\[2pt]
\quad FN-DSA-512  & 204.64 & 204.81 & 204.81 \\
\quad FN-DSA-1024 & 207.16 & 207.38 & 207.38 \\
\midrule
\multicolumn{4}{l}{\textit{SLH-DSA-SHAKE (FIPS 205)}} \\[2pt]
\quad SLH-DSA-SHAKE-128s & 210.31 & 210.31 & 210.31 \\
\quad SLH-DSA-SHAKE-128f & 219.32 & 219.32 & 219.32 \\
\quad SLH-DSA-SHAKE-192s & 218.52 & 218.52 & 218.52 \\
\quad SLH-DSA-SHAKE-192f & 237.51 & 237.51 & 237.51 \\
\quad SLH-DSA-SHAKE-256s & 231.82 & 231.82 & 231.82 \\
\quad SLH-DSA-SHAKE-256f & 251.42 & 251.42 & 251.42 \\
\bottomrule
\end{tabular}
\end{table}

%
% Pacotes necessários:
%   \usepackage{booktabs}
%   \usepackage{multirow}
%   \usepackage[utf8]{inputenc}
%
% Metodologia de medição:
%   - Stack  : stack painting via pthread (4 MB de buffer sentinela)
%   - Heap   : Valgrind --tool=massif, pico de mem_heap_B
% ============================================================


% ============================================================
% Tabela 1: Stack — KEMs
% ============================================================
\begin{table}[ht]
\centering
\caption{Uso de memória \textit{stack} por mecanismo de encapsulamento de
chave (KEM) e operação, em bytes. Medido via \textit{stack painting}
com \texttt{pthread} de 4\,MB e compilação \texttt{-O0}.}
\label{tab:stack_kem}
\begin{tabular}{lrrr}
\toprule
\textbf{Algoritmo} & \textbf{Keygen (B)} & \textbf{Encap (B)} & \textbf{Decap (B)} \\
\midrule
\multicolumn{4}{l}{\textit{ML-KEM (FIPS 203)}} \\[2pt]
\quad ML-KEM-512  &  13\,520 &  16\,160 &  16\,960 \\
\quad ML-KEM-768  &  17\,648 &  20\,784 &  21\,904 \\
\quad ML-KEM-1024 &  22\,768 &  26\,448 &  28\,048 \\
\midrule
\multicolumn{4}{l}{\textit{HQC}} \\[2pt]
\quad HQC-128 &  38\,576 &  56\,448 &  60\,896 \\
\quad HQC-192 &  70\,688 & 106\,752 & 115\,760 \\
\quad HQC-256 & 108\,768 & 166\,592 & 181\,040 \\
\bottomrule
\end{tabular}
\end{table}


% ============================================================
% Tabela 2: Stack — Assinaturas Digitais
% ============================================================
\begin{table}[ht]
\centering
\caption{Uso de memória \textit{stack} por esquema de assinatura digital
e operação, em bytes. Medido via \textit{stack painting}
com \texttt{pthread} de 4\,MB e compilação \texttt{-O0}.
Dados para SLH-DSA-SHAKE não disponíveis para esta métrica
(medição pendente com variantes SHAKE).}
\label{tab:stack_sig}
\begin{tabular}{lrrr}
\toprule
\textbf{Algoritmo} & \textbf{Keygen (B)} & \textbf{Sign (B)} & \textbf{Verify (B)} \\
\midrule
\multicolumn{4}{l}{\textit{ML-DSA (FIPS 204)}} \\[2pt]
\quad ML-DSA-44 &  28\,304 &  53\,456 &  25\,968 \\
\quad ML-DSA-65 &  33\,456 &  83\,312 &  29\,072 \\
\quad ML-DSA-87 &  41\,648 & 117\,072 &  35\,472 \\
\midrule
\multicolumn{4}{l}{\textit{FN-DSA (FIPS 206)}} \\[2pt]
\quad FN-DSA-512  &  23\,608 &  49\,584 &   9\,064 \\
\quad FN-DSA-1024 &  40\,504 &  90\,736 &  13\,160 \\
\midrule
\multicolumn{4}{l}{\textit{SLH-DSA-SHAKE (FIPS 205)}} \\[2pt]
\quad SLH-DSA-SHAKE-128s & --- & --- & --- \\
\quad SLH-DSA-SHAKE-128f & --- & --- & --- \\
\quad SLH-DSA-SHAKE-192s & --- & --- & --- \\
\quad SLH-DSA-SHAKE-192f & --- & --- & --- \\
\quad SLH-DSA-SHAKE-256s & --- & --- & --- \\
\quad SLH-DSA-SHAKE-256f & --- & --- & --- \\
\bottomrule
\end{tabular}
\end{table}


% ============================================================
% Tabela 3: Pico de Heap — KEMs (Valgrind massif)
% ============================================================
\begin{table}[ht]
\centering
\caption{Pico de uso de memória \textit{heap} por mecanismo de
encapsulamento de chave (KEM) e operação, em KB.
Medido com \texttt{valgrind --tool=massif}.
Os valores representam o máximo de \texttt{mem\_heap\_B}
ao longo da execução de cada operação isolada.}
\label{tab:heap_kem}
\begin{tabular}{lrrr}
\toprule
\textbf{Algoritmo} & \textbf{Keygen (KB)} & \textbf{Encap (KB)} & \textbf{Decap (KB)} \\
\midrule
\multicolumn{4}{l}{\textit{ML-KEM (FIPS 203)}} \\[2pt]
\quad ML-KEM-512  & 204.97 & 204.97 & 204.97 \\
\quad ML-KEM-768  & 207.17 & 207.17 & 207.17 \\
\quad ML-KEM-1024 & 208.00 & 208.17 & 208.17 \\
\midrule
\multicolumn{4}{l}{\textit{HQC}} \\[2pt]
\quad HQC-128 & 211.07 & 211.07 & 211.07 \\
\quad HQC-192 & 219.74 & 219.74 & 219.74 \\
\quad HQC-256 & 230.60 & 230.60 & 230.60 \\
\bottomrule
\end{tabular}
\end{table}


% ============================================================
% Tabela 4: Pico de Heap — Assinaturas Digitais (Valgrind massif)
% ============================================================
\begin{table}[ht]
\centering
\caption{Pico de uso de memória \textit{heap} por esquema de assinatura
digital e operação, em KB.
Medido com \texttt{valgrind --tool=massif}.
Os valores representam o máximo de \texttt{mem\_heap\_B}
ao longo da execução de cada operação isolada.}
\label{tab:heap_sig}
\begin{tabular}{lrrr}
\toprule
\textbf{Algoritmo} & \textbf{Keygen (KB)} & \textbf{Sign (KB)} & \textbf{Verify (KB)} \\
\midrule
\multicolumn{4}{l}{\textit{ML-DSA (FIPS 204)}} \\[2pt]
\quad ML-DSA-44 & 207.92 & 208.90 & 208.90 \\
\quad ML-DSA-65 & 211.02 & 211.02 & 211.83 \\
\quad ML-DSA-87 & 214.37 & 214.59 & 214.59 \\
\midrule
\multicolumn{4}{l}{\textit{FN-DSA (FIPS 206)}} \\[2pt]
\quad FN-DSA-512  & 204.64 & 204.81 & 204.81 \\
\quad FN-DSA-1024 & 207.16 & 207.38 & 207.38 \\
\midrule
\multicolumn{4}{l}{\textit{SLH-DSA-SHAKE (FIPS 205)}} \\[2pt]
\quad SLH-DSA-SHAKE-128s & 210.31 & 210.31 & 210.31 \\
\quad SLH-DSA-SHAKE-128f & 219.32 & 219.32 & 219.32 \\
\quad SLH-DSA-SHAKE-192s & 218.52 & 218.52 & 218.52 \\
\quad SLH-DSA-SHAKE-192f & 237.51 & 237.51 & 237.51 \\
\quad SLH-DSA-SHAKE-256s & 231.82 & 231.82 & 231.82 \\
\quad SLH-DSA-SHAKE-256f & 251.42 & 251.42 & 251.42 \\
\bottomrule
\end{tabular}
\end{table}

%
% Pacotes necessários:
%   \usepackage{booktabs}
%   \usepackage{multirow}
%   \usepackage[utf8]{inputenc}
%
% Metodologia de medição:
%   - Stack  : stack painting via pthread (4 MB de buffer sentinela)
%   - Heap   : Valgrind --tool=massif, pico de mem_heap_B
% ============================================================


% ============================================================
% Tabela 1: Stack — KEMs
% ============================================================
\begin{table}[ht]
\centering
\caption{Uso de memória \textit{stack} por mecanismo de encapsulamento de
chave (KEM) e operação, em bytes. Medido via \textit{stack painting}
com \texttt{pthread} de 4\,MB e compilação \texttt{-O0}.}
\label{tab:stack_kem}
\begin{tabular}{lrrr}
\toprule
\textbf{Algoritmo} & \textbf{Keygen (B)} & \textbf{Encap (B)} & \textbf{Decap (B)} \\
\midrule
\multicolumn{4}{l}{\textit{ML-KEM (FIPS 203)}} \\[2pt]
\quad ML-KEM-512  &  13\,520 &  16\,160 &  16\,960 \\
\quad ML-KEM-768  &  17\,648 &  20\,784 &  21\,904 \\
\quad ML-KEM-1024 &  22\,768 &  26\,448 &  28\,048 \\
\midrule
\multicolumn{4}{l}{\textit{HQC}} \\[2pt]
\quad HQC-128 &  38\,576 &  56\,448 &  60\,896 \\
\quad HQC-192 &  70\,688 & 106\,752 & 115\,760 \\
\quad HQC-256 & 108\,768 & 166\,592 & 181\,040 \\
\bottomrule
\end{tabular}
\end{table}


% ============================================================
% Tabela 2: Stack — Assinaturas Digitais
% ============================================================
\begin{table}[ht]
\centering
\caption{Uso de memória \textit{stack} por esquema de assinatura digital
e operação, em bytes. Medido via \textit{stack painting}
com \texttt{pthread} de 4\,MB e compilação \texttt{-O0}.
Dados para SLH-DSA-SHAKE não disponíveis para esta métrica
(medição pendente com variantes SHAKE).}
\label{tab:stack_sig}
\begin{tabular}{lrrr}
\toprule
\textbf{Algoritmo} & \textbf{Keygen (B)} & \textbf{Sign (B)} & \textbf{Verify (B)} \\
\midrule
\multicolumn{4}{l}{\textit{ML-DSA (FIPS 204)}} \\[2pt]
\quad ML-DSA-44 &  28\,304 &  53\,456 &  25\,968 \\
\quad ML-DSA-65 &  33\,456 &  83\,312 &  29\,072 \\
\quad ML-DSA-87 &  41\,648 & 117\,072 &  35\,472 \\
\midrule
\multicolumn{4}{l}{\textit{FN-DSA (FIPS 206)}} \\[2pt]
\quad FN-DSA-512  &  23\,608 &  49\,584 &   9\,064 \\
\quad FN-DSA-1024 &  40\,504 &  90\,736 &  13\,160 \\
\midrule
\multicolumn{4}{l}{\textit{SLH-DSA-SHAKE (FIPS 205)}} \\[2pt]
\quad SLH-DSA-SHAKE-128s & --- & --- & --- \\
\quad SLH-DSA-SHAKE-128f & --- & --- & --- \\
\quad SLH-DSA-SHAKE-192s & --- & --- & --- \\
\quad SLH-DSA-SHAKE-192f & --- & --- & --- \\
\quad SLH-DSA-SHAKE-256s & --- & --- & --- \\
\quad SLH-DSA-SHAKE-256f & --- & --- & --- \\
\bottomrule
\end{tabular}
\end{table}


% ============================================================
% Tabela 3: Pico de Heap — KEMs (Valgrind massif)
% ============================================================
\begin{table}[ht]
\centering
\caption{Pico de uso de memória \textit{heap} por mecanismo de
encapsulamento de chave (KEM) e operação, em KB.
Medido com \texttt{valgrind --tool=massif}.
Os valores representam o máximo de \texttt{mem\_heap\_B}
ao longo da execução de cada operação isolada.}
\label{tab:heap_kem}
\begin{tabular}{lrrr}
\toprule
\textbf{Algoritmo} & \textbf{Keygen (KB)} & \textbf{Encap (KB)} & \textbf{Decap (KB)} \\
\midrule
\multicolumn{4}{l}{\textit{ML-KEM (FIPS 203)}} \\[2pt]
\quad ML-KEM-512  & 204.97 & 204.97 & 204.97 \\
\quad ML-KEM-768  & 207.17 & 207.17 & 207.17 \\
\quad ML-KEM-1024 & 208.00 & 208.17 & 208.17 \\
\midrule
\multicolumn{4}{l}{\textit{HQC}} \\[2pt]
\quad HQC-128 & 211.07 & 211.07 & 211.07 \\
\quad HQC-192 & 219.74 & 219.74 & 219.74 \\
\quad HQC-256 & 230.60 & 230.60 & 230.60 \\
\bottomrule
\end{tabular}
\end{table}


% ============================================================
% Tabela 4: Pico de Heap — Assinaturas Digitais (Valgrind massif)
% ============================================================
\begin{table}[ht]
\centering
\caption{Pico de uso de memória \textit{heap} por esquema de assinatura
digital e operação, em KB.
Medido com \texttt{valgrind --tool=massif}.
Os valores representam o máximo de \texttt{mem\_heap\_B}
ao longo da execução de cada operação isolada.}
\label{tab:heap_sig}
\begin{tabular}{lrrr}
\toprule
\textbf{Algoritmo} & \textbf{Keygen (KB)} & \textbf{Sign (KB)} & \textbf{Verify (KB)} \\
\midrule
\multicolumn{4}{l}{\textit{ML-DSA (FIPS 204)}} \\[2pt]
\quad ML-DSA-44 & 207.92 & 208.90 & 208.90 \\
\quad ML-DSA-65 & 211.02 & 211.02 & 211.83 \\
\quad ML-DSA-87 & 214.37 & 214.59 & 214.59 \\
\midrule
\multicolumn{4}{l}{\textit{FN-DSA (FIPS 206)}} \\[2pt]
\quad FN-DSA-512  & 204.64 & 204.81 & 204.81 \\
\quad FN-DSA-1024 & 207.16 & 207.38 & 207.38 \\
\midrule
\multicolumn{4}{l}{\textit{SLH-DSA-SHAKE (FIPS 205)}} \\[2pt]
\quad SLH-DSA-SHAKE-128s & 210.31 & 210.31 & 210.31 \\
\quad SLH-DSA-SHAKE-128f & 219.32 & 219.32 & 219.32 \\
\quad SLH-DSA-SHAKE-192s & 218.52 & 218.52 & 218.52 \\
\quad SLH-DSA-SHAKE-192f & 237.51 & 237.51 & 237.51 \\
\quad SLH-DSA-SHAKE-256s & 231.82 & 231.82 & 231.82 \\
\quad SLH-DSA-SHAKE-256f & 251.42 & 251.42 & 251.42 \\
\bottomrule
\end{tabular}
\end{table}
